\chapter{Modello dei dati}
\label{chap:database}

\section{Schema relazionale}
\label{sec:schema}

Lo schema relazionale (definito in \texttt{init.sql}) si articola in nove tabelle principali.
Tutte le chiavi primarie sono \texttt{UUID} generati dall'estensione \texttt{pgcrypto}
(\texttt{gen\_random\_uuid()}). Le foreign key utilizzano \texttt{ON DELETE CASCADE}
ovunque la relazione sia di composizione, garantendo l'invariante di assenza di
\emph{dangling references}.

\begin{figure}[ht]
\footnotesize
\begin{verbatim}
users  ─── 1:N ─── folders        (gerarchia con parent_id self-ref)
   │                  │
   │                  └── 1:N ─── documents
   │
   ├── 1:N ─── permissions        (target esclusivo: folder XOR document)
   ├── 1:N ─── history
   ├── 1:N ─── email_accounts
   │              │
   │              └── 1:N ─── email_jobs ─── 1:N ─── agent_files
   │                                │                      │
   │                                └── 1:N ─── agent_operations
   │
   └── 1:N ─── agent_operations
\end{verbatim}
\caption{Sintesi testuale del modello entit\`a-relazione.}
\label{fig:erd}
\end{figure}

\section{Tabelle principali}
\label{sec:tabelle_principali}

\subsection{users}

Tabella centrale del sistema: rappresenta un account utente.

\begin{lstlisting}[language=SQL, caption={DDL della tabella \texttt{users}.}]
CREATE TABLE users (
    id            UUID PRIMARY KEY DEFAULT gen_random_uuid(),
    username      VARCHAR(100) NOT NULL,
    email         VARCHAR(100) UNIQUE NOT NULL,
    password_hash TEXT NOT NULL,
    role          VARCHAR(20) DEFAULT 'USER',   -- USER | DOMAIN_ADMIN
    is_active     BOOLEAN DEFAULT TRUE,
    created_at    TIMESTAMP WITH TIME ZONE DEFAULT now(),
    last_active_at TIMESTAMP WITH TIME ZONE DEFAULT now()
);
\end{lstlisting}

Il campo \texttt{last\_active\_at} \`e aggiornato dal middleware \texttt{LastActiveMiddleware}
con debounce di 5~minuti per ogni richiesta autenticata, evitando scritture eccessive
per ogni singola API call.

\subsection{folders}

Struttura gerarchica ad albero realizzata tramite self-join su \texttt{parent\_id}:

\begin{lstlisting}[language=SQL, caption={DDL della tabella \texttt{folders}.}]
CREATE TABLE folders (
    id         UUID PRIMARY KEY DEFAULT gen_random_uuid(),
    name       VARCHAR(255) NOT NULL,
    parent_id  UUID REFERENCES folders(id) ON DELETE CASCADE,
    owner_id   UUID NOT NULL REFERENCES users(id) ON DELETE CASCADE,
    created_at TIMESTAMP WITH TIME ZONE DEFAULT now()
);
CREATE INDEX idx_folders_parent_id ON folders(parent_id);
\end{lstlisting}

\subsection{documents}

\begin{lstlisting}[language=SQL, caption={DDL della tabella \texttt{documents}.}]
CREATE TABLE documents (
    id         UUID PRIMARY KEY DEFAULT gen_random_uuid(),
    name       VARCHAR(255) NOT NULL,
    mime_type  VARCHAR(100),
    size_bytes BIGINT,           -- NULL fino alla conferma post-upload
    folder_id  UUID REFERENCES folders(id) ON DELETE CASCADE,
    owner_id   UUID NOT NULL REFERENCES users(id) ON DELETE CASCADE,
    created_at TIMESTAMP WITH TIME ZONE DEFAULT now(),
    updated_at TIMESTAMP WITH TIME ZONE DEFAULT now()
);
CREATE INDEX idx_documents_folder_id ON documents(folder_id);
CREATE INDEX idx_documents_owner_id  ON documents(owner_id);
\end{lstlisting}

Il campo \texttt{size\_bytes = NULL} indica un documento in attesa di conferma
(fase 1 del two-phase upload): \`e filtrato dalle listing degli utenti.

\subsection{permissions}

\begin{lstlisting}[language=SQL, caption={DDL della tabella \texttt{permissions}.}]
CREATE TABLE permissions (
    id           UUID PRIMARY KEY DEFAULT gen_random_uuid(),
    user_id      UUID NOT NULL REFERENCES users(id) ON DELETE CASCADE,
    folder_id    UUID REFERENCES folders(id) ON DELETE CASCADE,
    document_id  UUID REFERENCES documents(id) ON DELETE CASCADE,
    access_level VARCHAR(20) NOT NULL,
    shared_at    TIMESTAMP WITH TIME ZONE DEFAULT now(),
    CONSTRAINT chk_permission_target CHECK (
        (folder_id IS NOT NULL AND document_id IS NULL) OR
        (folder_id IS NULL AND document_id IS NOT NULL)
    ),
    CONSTRAINT chk_access_level CHECK (
        access_level IN ('VIEWER', 'EDITOR')
    )
);
\end{lstlisting}

Il vincolo \texttt{chk\_permission\_target} garantisce che ogni permesso si riferisca
a esattamente una risorsa (folder XOR document). La cascata di delete dalla risorsa
padre revoca automaticamente i permessi su di essa.

\section{Tabelle per l'integrazione email}
\label{sec:tabelle_email}

\subsection{email\_accounts}

\begin{lstlisting}[language=SQL, caption={DDL della tabella \texttt{email\_accounts} (estratto).}]
CREATE TABLE email_accounts (
    id                    UUID PRIMARY KEY DEFAULT gen_random_uuid(),
    user_id               UUID NOT NULL REFERENCES users(id) ON DELETE CASCADE,
    email_address         VARCHAR(255) NOT NULL,
    imap_host             VARCHAR(255) NOT NULL,
    imap_port             INTEGER DEFAULT 993,
    use_ssl               BOOLEAN DEFAULT TRUE,
    auth_type             VARCHAR(20) DEFAULT 'app_password',
    encrypted_credentials TEXT NOT NULL,   -- Fernet-encrypted
    oauth_provider        VARCHAR(20),
    is_active             BOOLEAN DEFAULT TRUE,
    last_synced_at        TIMESTAMP WITH TIME ZONE,
    last_uid              BIGINT DEFAULT 0,  -- watermark IMAP
    created_at            TIMESTAMP WITH TIME ZONE DEFAULT now(),
    CONSTRAINT chk_auth_type CHECK (auth_type IN ('app_password', 'oauth2'))
);
\end{lstlisting}

Il campo \texttt{last\_uid} \`e il watermark IMAP: ad ogni ciclo il poller recupera
solo i messaggi con UID superiore a questo valore, garantendo l'idempotenza rispetto
ai messaggi gi\`a processati.

\subsection{email\_jobs}

\begin{lstlisting}[language=SQL, caption={DDL della tabella \texttt{email\_jobs}.}]
CREATE TABLE email_jobs (
    id               UUID PRIMARY KEY DEFAULT gen_random_uuid(),
    email_account_id UUID NOT NULL REFERENCES email_accounts(id) ON DELETE CASCADE,
    message_uid      BIGINT NOT NULL,
    subject          TEXT,
    sender           VARCHAR(255),
    received_at      TIMESTAMP WITH TIME ZONE,
    eml_storage_key  TEXT,   -- chiave S3 del file .eml grezzo
    status           VARCHAR(20) DEFAULT 'pending',
    error_message    TEXT,
    created_at       TIMESTAMP WITH TIME ZONE DEFAULT now(),
    processed_at     TIMESTAMP WITH TIME ZONE,
    CONSTRAINT chk_job_status CHECK (
        status IN ('pending', 'processing', 'done', 'failed')
    ),
    UNIQUE (email_account_id, message_uid)   -- idempotenza IMAP
);
\end{lstlisting}

Il vincolo \texttt{UNIQUE(email\_account\_id, message\_uid)} garantisce che lo stesso
messaggio non venga processato due volte, anche in presenza di riavvii del poller o di
esecuzioni concorrenti.

\section{Tabelle per il layer agentico}
\label{sec:tabelle_agentiche}

\subsection{agent\_files}

\begin{lstlisting}[language=SQL, caption={DDL della tabella \texttt{agent\_files}.}]
CREATE TABLE agent_files (
    id                 UUID PRIMARY KEY DEFAULT gen_random_uuid(),
    job_id             UUID REFERENCES email_jobs(id) ON DELETE CASCADE,
    source             VARCHAR(20) NOT NULL DEFAULT 'email',
    filename           VARCHAR(255) NOT NULL,
    mime_type          VARCHAR(100),
    size_bytes         BIGINT,
    content_hash       VARCHAR(64),   -- SHA-256 per deduplicazione
    suggested_folder_id UUID REFERENCES folders(id) ON DELETE SET NULL,
    confidence         FLOAT,
    agent_reasoning    TEXT,
    document_id        UUID REFERENCES documents(id) ON DELETE SET NULL,
    status             VARCHAR(20) DEFAULT 'pending',
    auto_filed         BOOLEAN DEFAULT FALSE,
    needs_classification BOOLEAN NOT NULL DEFAULT TRUE,
    created_at         TIMESTAMP WITH TIME ZONE DEFAULT now(),
    CONSTRAINT chk_agent_file_status CHECK (
        status IN ('pending','auto_filed','in_inbox','confirmed',
                   'rejected','indexed','failed')
    ),
    CONSTRAINT chk_agent_file_source CHECK (
        source IN ('email', 'manual_upload')
    ),
    CONSTRAINT chk_agent_file_source_has_job CHECK (
        (source = 'email' AND job_id IS NOT NULL) OR
        (source = 'manual_upload' AND job_id IS NULL)
    )
);
\end{lstlisting}

Il campo \texttt{content\_hash} (SHA-256 del contenuto del file) \`e usato per la
deduplicazione: file identici ricevuti pi\`u volte non vengono ri-classificati.
Il flag \texttt{needs\_classification} separa i caricamenti manuali diretti (che
richiedono solo indicizzazione in OpenSearch) dai caricamenti AI-assisted (che
richiedono anche l'invocazione dell'agente LLM).

\subsection{agent\_operations}

\begin{lstlisting}[language=SQL, caption={DDL della tabella \texttt{agent\_operations}.}]
CREATE TABLE agent_operations (
    id             UUID PRIMARY KEY DEFAULT gen_random_uuid(),
    user_id        UUID NOT NULL REFERENCES users(id) ON DELETE CASCADE,
    job_id         UUID REFERENCES email_jobs(id),
    agent_file_id  UUID REFERENCES agent_files(id),
    operation_type VARCHAR(50) NOT NULL,
    description    TEXT NOT NULL,
    details        JSONB,           -- struttura arbitraria per evento
    created_at     TIMESTAMP WITH TIME ZONE DEFAULT now(),
    CONSTRAINT chk_operation_type CHECK (
        operation_type IN (
            'email_received', 'attachment_extracted', 'auto_filed',
            'sent_to_inbox', 'duplicate_skipped',
            'manual_upload_received', 'error'
        )
    )
);
CREATE INDEX idx_agent_operations_user_id    ON agent_operations(user_id);
CREATE INDEX idx_agent_operations_created_at ON agent_operations(created_at);
\end{lstlisting}

Questa tabella costituisce l'audit log immutabile delle azioni agentiche, esposto
all'endpoint \texttt{GET /agent/operations} e visualizzato nella \emph{Agent History
view}. La colonna \texttt{details JSONB} accoglie metadati strutturati specifici per
tipo di evento (es.\ \texttt{\{"folder\_id": "...", "confidence": 0.92\}} per
\texttt{auto\_filed}).

\section{Macchine a stati}
\label{sec:state_machines}

\subsection{email\_jobs}

\begin{table}[ht]
\centering
\small
\begin{tabularx}{\linewidth}{l X l}
\toprule
\textbf{Stato} & \textbf{Significato} & \textbf{Transizioni in uscita} \\
\midrule
\texttt{pending}    & Job creato dal poller, nessun lavoro avviato.                  & $\rightarrow$ \texttt{processing} \\
\texttt{processing} & Worker ha preso in carico il job.                              & $\rightarrow$ \texttt{done} / \texttt{failed} \\
\texttt{done}       & Allegati estratti, classificati e persistiti.                  & --- \\
\texttt{failed}     & Errore non recuperabile; \texttt{error\_message} valorizzato.  & manuale \\
\bottomrule
\end{tabularx}
\caption{Macchina a stati di \texttt{email\_jobs}.}
\label{tab:state_machine_jobs}
\end{table}

\subsection{agent\_files}

\begin{table}[ht]
\centering
\small
\begin{tabularx}{\linewidth}{l X}
\toprule
\textbf{Stato} & \textbf{Significato} \\
\midrule
\texttt{pending}      & File creato, in attesa di elaborazione da parte dell'agent worker. \\
\texttt{auto\_filed}  & Classificato con confidenza $\geq \theta$; documento spostato nella cartella. \\
\texttt{in\_inbox}    & Confidenza $< \theta$ o nessuna cartella valida; proposta in attesa di revisione. \\
\texttt{indexed}      & Solo indicizzazione OpenSearch richiesta (no classificazione); completato. \\
\texttt{confirmed}    & Proposta accettata dall'utente dall'\emph{inbox view}. \\
\texttt{rejected}     & Proposta rigettata dall'utente. \\
\texttt{failed}       & Errore non recuperabile (Tika irraggiungibile, LLM error, ecc.). \\
\bottomrule
\end{tabularx}
\caption{Macchina a stati di \texttt{agent\_files}.}
\label{tab:state_machine_files}
\end{table}

\section{Strategia di indicizzazione}
\label{sec:indexing}

Gli indici sono definiti sulle colonne usate pi\`u frequentemente nelle query critiche:

\begin{itemize}[leftmargin=*]
    \item \texttt{folders(parent\_id)}: navigazione gerarchica dell'albero cartelle.
    \item \texttt{documents(folder\_id)}, \texttt{documents(owner\_id)}: listing documenti per cartella o proprietario.
    \item \texttt{permissions(user\_id)}, \texttt{permissions(user\_id, folder\_id)},
          \texttt{permissions(user\_id, document\_id)}: verifica dei permessi su ogni operazione.
    \item \texttt{email\_jobs(status)}: selezione dei job \texttt{pending} da parte del worker.
    \item \texttt{agent\_files(status)}, \texttt{agent\_files(source, status)}: selezione dei
          file da elaborare nei due rami del worker (email vs.\ manual\_upload).
    \item \texttt{agent\_files(content\_hash)}: deduplicazione rapida.
    \item \texttt{agent\_operations(user\_id)}, \texttt{agent\_operations(created\_at)}:
          timeline delle operazioni agentiche.
\end{itemize}

\section{Tabella di branding multi-tenant}
\label{sec:branding_table}

\begin{lstlisting}[language=SQL, caption={DDL della tabella \texttt{domain\_branding}.}]
CREATE TABLE domain_branding (
    domain          VARCHAR(255) PRIMARY KEY,  -- e.g. "azienda.com"
    brand_name      VARCHAR(100),
    primary_color   VARCHAR(20),   -- colore hex
    logo            BYTEA,         -- immagine binaria
    logo_mime_type  VARCHAR(100),
    updated_at      TIMESTAMP WITH TIME ZONE DEFAULT now()
);
\end{lstlisting}

Ogni tenant pu\`o personalizzare il nome del brand, il colore primario dell'interfaccia
e il logo aziendale. La chiave primaria \`e il dominio email (es.\ \texttt{azienda.com}),
derivato automaticamente alla registrazione del primo utente.
